% ============================================================================
%  Actividad 4 - El tema institucional y los cinco flujos de tarea
%
%  Propiedad de US-ENTREGA-A4. Tres bloques y nada mas: el tema institucional
%  con la procedencia de cada valor, su tipografia y su matriz de contraste en
%  los dos modos; los cinco flujos de tarea del archivo de diseno; y la
%  verificacion tecnica de ese archivo.
%
%  Regla de este archivo: ningun valor de color se escribe con almohadilla.
%  scripts/verificar_tokens_a4.sh falla si encuentra un hexadecimal con
%  almohadilla en cualquier contenido/a4_*.tex. Los valores se publican como
%  texto monoespaciado, que es lo que son aqui: el dato de procedencia de un
%  token del portal, no una instruccion de color de este informe.
%
%  Regla de no duplicacion: a4_03_guia_estilos.tex ya cubre identidad,
%  tipografia del tema de omision, reticula, microinteracciones, iconografia,
%  imagenes, voz y tono y bitacora de versiones. Aqui no se repite ninguna.
% ============================================================================

% ============================================================================
\uxsection{El tema institucional}

La guía de estilos de este documento describe el sistema de diseño con el que están construidas
las siete pantallas del portal y tomadas sus capturas. Ese sistema entrega además un segundo tema,
derivado del archivo de diseño de esta actividad, y esta sección lo documenta con el mismo rigor:
de dónde sale cada valor, qué familia tipográfica lo acompaña y qué razón de contraste alcanza cada
token en los dos modos de color.

Los dos temas conviven y ninguno sustituye al otro. El tema de omisión, que es el que las capturas
de este documento muestran, conserva sus valores sin una sola diferencia. El tema institucional se
elige desde la cabecera del portal, viaja en una cookie y el servidor lo aplica antes del primer
trazado, de modo que el cambio no produce destello. Cada tema ofrece modo claro y modo oscuro:
\textbf{cuatro combinaciones}, verificadas una por una con el generador que emite los tokens.

\subsection{El octeto del archivo de diseño}

El archivo de diseño declara ocho colores con su función escrita al lado. Son el origen de la
paleta institucional y la referencia contra la que se mide cualquier derivación.

\begin{uxtabla}{|L{2.6cm}|C{2.0cm}|L{3.5cm}|Y|}{Los ocho colores que declara el archivo de diseño y su destino en el sistema}
  \uxheadrow \thd{Nombre en el archivo} & \thd{Valor} & \thd{Función declarada} & \thd{Destino en el contrato de tokens} \\
  Navegación & \texttt{102A43} & Estructura de navegación & Corriente plena en claro y origen del suelo oscuro \\
  Secundario & \texttt{1D4C6E} & Segundo nivel de estructura & Corriente media en claro \\
  Acción & \texttt{086B70} & Acciones y selección & Ninguno: se aplica desde la capa de componentes \\
  Apoyo & \texttt{15989A} & Acento de apoyo & Ninguno: con 3{,}51:1 sobre la superficie no puede llevar texto \\
  Atención & \texttt{B97812} & Atención & Aviso, con la corrección de luminancia documentada abajo \\
  Éxito & \texttt{287A58} & Confirmación & Confirmación, con su valor literal \\
  Error & \texttt{B8443F} & Error & Error, con la corrección de luminancia documentada abajo \\
  Superficie & \texttt{FFFFFF} & Superficie de trabajo & Suelo del modo claro, con su valor literal \\
\end{uxtabla}

Seis de los ocho ocupan una ranura del contrato de tokens; dos no la ocupan, y en los dos casos por
una razón escrita. La consecuencia práctica es que el tema institucional no se limita a copiar el
archivo: toma de él lo que el contrato puede alojar, deriva el resto con una regla declarada y mide
el resultado antes de publicarlo.

\subsection{Los diecisiete tokens, con la procedencia de cada valor}

El contrato del sistema tiene diecisiete tokens de color y el tema institucional los resuelve todos
en los dos modos. La columna de procedencia distingue tres orígenes posibles: el valor literal del
archivo de diseño, una derivación declarada a partir de uno de sus colores, o el valor heredado del
tema de omisión.

\begin{uxtablalarga}{|L{3.4cm}|C{1.9cm}|C{1.9cm}|Y|}{La paleta institucional: diecisiete tokens en los dos modos, con la procedencia de cada valor}{\thd{Token} & \thd{Claro} & \thd{Oscuro} & \thd{Procedencia}}
  \texttt{ground} & \texttt{FFFFFF} & \texttt{0B1B2B} & Claro: superficie del archivo. Oscuro: derivado del color de navegación bajando su luminancia; azul profundo y nunca negro puro \\
  \texttt{ground-alt} & \texttt{F1F4F8} & \texttt{102A43} & Claro: derivado del color de navegación, aclarado hasta 1{,}10:1 sobre la superficie. Oscuro: el color de navegación usado como panel \\
  \texttt{grid} & \texttt{DCE3EC} & \texttt{1D3348} & Derivado en los dos modos: un paso del suelo hacia la rampa. Decorativo por definición \\
  \texttt{corriente-apagado} & \texttt{A8B4C4} & \texttt{46586B} & Derivado. El peldaño se queda a propósito por debajo de 3:1, porque no informa \\
  \texttt{corriente-tenue} & \texttt{4A5A6E} & \texttt{93A6BC} & Derivado del color de navegación hacia el suelo hasta cruzar 4{,}5:1 \\
  \texttt{corriente-medio} & \texttt{1D4C6E} & \texttt{BFD0E2} & Claro: color secundario del archivo. Oscuro: el mismo matiz elevado sobre el suelo profundo \\
  \texttt{corriente-pleno} & \texttt{102A43} & \texttt{EAF2FA} & Claro: color de navegación del archivo. Oscuro: el mismo matiz elevado \\
  \texttt{error} & \texttt{933632} & \texttt{EC5B51} & Derivado del color de error del archivo, oscurecido conservando matiz y saturación \\
  \texttt{aviso} & \texttt{A36A10} & \texttt{E8A33D} & Claro: color de atención oscurecido un 12 \%. Oscuro: el mismo color aclarado para el suelo profundo \\
  \texttt{ok} & \texttt{287A58} & \texttt{5FCB94} & Claro: color de éxito del archivo. Oscuro: el mismo color aclarado \\
  \texttt{info} & \texttt{123C7A} & \texttt{B9A7F2} & Derivado y deliberadamente fuera del octeto, por la razón que la subsección de hallazgos documenta \\
  \texttt{serie-1} & \texttt{1D4ED8} & \texttt{7DD3FC} & Heredado del tema de omisión \\
  \texttt{serie-2} & \texttt{B45309} & \texttt{FFC233} & Heredado del tema de omisión \\
  \texttt{serie-3} & \texttt{1F6F43} & \texttt{4ADE80} & Heredado del tema de omisión \\
  \texttt{serie-4} & \texttt{6D28D9} & \texttt{C4B5FD} & Heredado del tema de omisión \\
  \texttt{serie-5} & \texttt{0E7490} & \texttt{67E8F9} & Heredado del tema de omisión \\
  \texttt{serie-6} & \texttt{9D174D} & \texttt{F9A8D4} & Heredado del tema de omisión \\
\end{uxtablalarga}

Los seis colores de serie \textbf{no cambian con el tema}, y es una decisión y no un olvido: una
serie es un canal de datos y no identidad de marca, el archivo de diseño no declara paleta
categórica y las seis ya están verificadas bajo tres dicromacias. Lo que sí se vuelve a medir por
tema es su razón sobre cada suelo, porque el suelo oscuro institucional es \texttt{0B1B2B} y no el
del tema de omisión: las seis quedan por encima de 3:1 en las cuatro combinaciones, con un mínimo
de \textbf{5{,}02:1} en claro y \textbf{9{,}43:1} en oscuro.

\subsection{La familia tipográfica del tema}

El eje del tema son dos cosas: el color y la familia tipográfica. El archivo de diseño declara la
suya con su propio criterio: \enquote{Inter mantiene claridad en tablas, filtros y metadatos. Los
tamaños priorizan densidad sin reducir la legibilidad; el espaciado entre letras permanece en 0.}

Bajo el tema institucional, los roles de titular y de cuerpo se componen con \textbf{Inter},
servida por el mismo módulo de fuentes y con el mismo proveedor que las dos familias del tema de
omisión. La familia monoespaciada \textbf{no cambia}: el rol de dato existe por el ancho fijo de
las cifras y no por identidad, y el archivo de diseño no declara ninguna monoespaciada. Los nueve
roles conservan tamaño, interlínea y espaciado entre letras, de modo que el tema cambia la letra y
no la jerarquía.

\figuraux{figuras/a4/guia_tipografia.png}{Lámina de tipografía del archivo de diseño, con el
criterio que declara la familia y su aplicación a títulos, métricas, contenido y controles. Es
diseño de alta fidelidad y no una captura del portal en ejecución.}{fig:a4-tema-tipografia}

\subsection{Matriz de contraste del tema institucional en los dos modos}

Las razones siguientes no se estimaron: las calcula el mismo generador que emite los tokens del
portal, con la fórmula de luminancia relativa de la norma y contra sus umbrales de 4{,}5:1 para
texto normal, 3:1 para elemento gráfico y 7:1 para el nivel más alto (W3C, 2023). Cada token se
mide contra el suelo de su modo: \texttt{FFFFFF} en claro y \texttt{0B1B2B} en oscuro.

\begin{uxtablalarga}{|L{3.4cm}|C{2.1cm}|C{2.0cm}|C{2.1cm}|Y|}{Matriz de contraste del tema institucional sobre el suelo de cada modo}{\thd{Token} & \thd{Claro} & \thd{Veredicto} & \thd{Oscuro} & \thd{Veredicto}}
  \texttt{ground-alt} & 1{,}10:1 & superficie & 1{,}19:1 & superficie \\
  \texttt{grid} & 1{,}29:1 & decorativo & 1{,}34:1 & decorativo \\
  \texttt{corriente-apagado} & 2{,}10:1 & decorativo & 2{,}38:1 & decorativo \\
  \texttt{corriente-tenue} & 7{,}05:1 & AAA & 6{,}98:1 & AA \\
  \texttt{corriente-medio} & 9{,}09:1 & AAA & 11{,}06:1 & AAA \\
  \texttt{corriente-pleno} & 14{,}64:1 & AAA & 15{,}41:1 & AAA \\
  \texttt{error} & 7{,}46:1 & AAA & 5{,}13:1 & AA \\
  \texttt{aviso} & 4{,}54:1 & AA & 8{,}07:1 & AAA \\
  \texttt{ok} & 5{,}23:1 & AA & 8{,}67:1 & AAA \\
  \texttt{info} & 10{,}75:1 & AAA & 8{,}19:1 & AAA \\
  \texttt{serie-1} & 6{,}70:1 & AA & 10{,}44:1 & AAA \\
  \texttt{serie-2} & 5{,}02:1 & AA & 10{,}79:1 & AAA \\
  \texttt{serie-3} & 6{,}15:1 & AA & 9{,}99:1 & AAA \\
  \texttt{serie-4} & 7{,}10:1 & AAA & 9{,}43:1 & AAA \\
  \texttt{serie-5} & 5{,}36:1 & AA & 12{,}01:1 & AAA \\
  \texttt{serie-6} & 7{,}88:1 & AAA & 9{,}60:1 & AAA \\
\end{uxtablalarga}

Dos veredictos de la tabla no son grados de la norma sino papeles del sistema. \textbf{Superficie}
marca un fondo medido contra otro fondo: nada se lee nunca con la fila alterna como primer plano,
así que exigirle 4{,}5:1 sería acusar al sistema de un incumplimiento que no existe.
\textbf{Decorativo} marca los dos tokens que declaran no informar, la retícula y el conector en
reposo, y su exención es la razón misma por la que esa declaración existe en el contrato.

El resultado agregado es el que importa: \textbf{cero incumplimientos} en las cuatro combinaciones
de tema y modo. Ningún token que informa se queda por debajo de 3:1 y ningún token que declara no
informar lo alcanza, que son las dos formas de romper la regla.

\subsection{Separación bajo dicromacia, medida por pares}

El contraste sobre el suelo no basta para el canal semántico: dos marcas pueden cumplir el umbral
por separado y confundirse entre sí para quien tiene una dicromacia. Por eso los cuatro semánticos
se miden también dos a dos, simulando protanopia, deuteranopia y tritanopia, y de cada par se
publica la distancia \textbf{peor} de las tres.

\begin{uxtabla}{|L{2.8cm}|C{1.8cm}|L{2.6cm}|C{1.8cm}|Y|}{Separación de los pares semánticos del tema institucional, con la dicromacia que produce la peor distancia}
  \uxheadrow \thd{Par} & \thd{Claro} & \thd{Dicromacia} & \thd{Oscuro} & \thd{Dicromacia} \\
  Error y aviso & 20{,}7 & tritanopia & 22{,}6 & deuteranopia \\
  Error y confirmación & 14{,}5 & protanopia & 23{,}3 & protanopia \\
  Error e informativo & 62{,}6 & protanopia & 67{,}6 & tritanopia \\
  Aviso y confirmación & 37{,}8 & protanopia & 37{,}8 & protanopia \\
  Aviso e informativo & 59{,}9 & tritanopia & 35{,}3 & tritanopia \\
  Confirmación e informativo & 20{,}1 & tritanopia & 25{,}9 & tritanopia \\
\end{uxtabla}

El piso no es un número inventado para esta paleta: es la peor separación que el tema de omisión ya
sostiene, \textbf{13{,}6} en claro y \textbf{21{,}5} en oscuro. El tema institucional queda por
encima de los dos, con \textbf{14{,}5} y \textbf{22{,}6}. Un tema nuevo que empeorara ese número
habría degradado el sistema aunque cada token cumpliera su umbral por separado.

\subsection{Los tres hallazgos de la verificación}

El cálculo no confirmó la paleta del archivo: la corrigió en tres puntos, y los tres se corrigieron
en el valor del token y no en la prosa.

\begin{enumerate}[leftmargin=*,itemsep=5pt,topsep=6pt]
\item \textbf{El color de atención no alcanzaba el listón de texto.} El valor del archivo,
  \texttt{B97812}, da \textbf{3{,}65:1} sobre la superficie blanca, por debajo del 4{,}5:1 que exige
  un token que informa. Oscurecido un 12 \% conservando matiz y saturación llega a \texttt{A36A10}
  con \textbf{4{,}54:1}. Es el mínimo que cruza: a un 10 \% todavía da 4{,}40:1.
\item \textbf{El canal informativo necesita un matiz fuera del octeto.} El color de acción y el de
  éxito se separan por \textbf{6{,}7} bajo tritanopia, muy por debajo del piso de 13{,}6: para
  quien no distingue el azul del amarillo, una marca informativa y una confirmación se ven igual.
  En claro se resuelve con un azul profundo derivado, \texttt{123C7A}, que lleva el peor par de ese
  canal a \textbf{20{,}1}. En oscuro no hay azul que sirva, porque ahí el azul y el verde colapsan,
  y la solución es un violeta, \texttt{B9A7F2}, con \textbf{25{,}9} sobre un piso de 21{,}5.
\item \textbf{El color de error del archivo no podía embarcarse tal cual.} Con \texttt{B8443F} el
  error da 5{,}33:1 y la confirmación 5{,}23:1, de modo que bajo protanopia los dos se separan por
  \textbf{10{,}1} contra un piso de 13{,}6 y un error se lee como una confirmación. Oscurecido un
  20 \%, \texttt{933632} sube a 7{,}46:1 y su peor par a \textbf{14{,}5}. En oscuro, el mismo color
  aclarado para el suelo profundo se quedaba en \textbf{13{,}1} contra un piso de 21{,}5;
  oscurecido un 12 \%, \texttt{EC5B51} da 5{,}13:1 y su peor par sube a \textbf{22{,}6}.
\end{enumerate}

Los tres siguen la misma regla, escrita antes de aplicarla: se corrige la luminancia y se conservan
el matiz y la saturación, de manera que el color siga siendo reconocible como el del archivo. Y los
tres se descubrieron por cálculo y no a ojo, que es la única forma de encontrar un defecto que solo
existe para un lector con dicromacia.

\subsection{El color de acción no entra en los diecisiete tokens}

El archivo declara un color de acción, \texttt{086B70}, y el sistema no le abre una ranura. La
razón es de contrato y no de gusto: los diecisiete tokens cubren suelo, retícula, rampa de
corriente, semántica y series, y ninguno de esos papeles es \enquote{acción}. Añadir un token
significa cambiar el contrato que consumen la aplicación, la paleta tipada y este documento, y ese
trabajo no cabía en esta entrega.

Mientras tanto el color de acción se aplica \textbf{desde la capa de componentes}, que es donde el
archivo de diseño lo usa: botón primario, selección y foco. La decisión queda declarada aquí en vez
de resolverse en silencio, porque un color que vive fuera del contrato es exactamente el tipo de
valor que se duplica a mano en tres componentes y se separa de su origen sin que nadie lo note.

\subsection{Los ocho estados interactivos, declarados como derivación}

El archivo de diseño publica ocho estados interactivos y los declara por \textbf{nombre de
variable}, no por valor: la lámina muestra el color y nombra la variable que lo produce. Esta
entrega no inventa los ocho valores ni los copia de la lámina con un cuentagotas: los deriva con
una regla escrita y marca la derivación como tal.

\textbf{La regla.} Sobre el color base, un \textbf{12 \%} de oscurecimiento para el estado al pasar
el puntero y un \textbf{20 \%} para el estado presionado, conservando matiz y saturación. Es la
misma operación con la que se corrigieron los tres hallazgos de la subsección anterior, de modo que
el sistema tiene una sola forma de mover un color.

\begin{uxtabla}{|L{4.2cm}|L{3.4cm}|Y|}{Los ocho estados interactivos del archivo de diseño y su procedencia}
  \uxheadrow \thd{Estado} & \thd{Base} & \thd{Procedencia} \\
  Primario al pasar & Color de acción & Derivación: base oscurecida un 12 \% \\
  Primario presionado & Color de acción & Derivación: base oscurecida un 20 \% \\
  Secundario al pasar & Color de navegación & Derivación: base oscurecida un 12 \% \\
  Secundario presionado & Color de navegación & Derivación: base oscurecida un 20 \% \\
  Deshabilitado & Neutro de la rampa & Derivación: peldaño apagado, sin papel informativo \\
  Foco & Color de acción & Variable propia del archivo, aplicada como filete de foco \\
  Error & Color de error & Variable propia del archivo, con el valor corregido del token \\
  Información & Canal informativo & Variable propia del archivo, con el valor derivado del token \\
\end{uxtabla}

\figuraux{figuras/a4/guia_paleta.png}{Lámina de paleta y estados interactivos del archivo de
diseño. Los ocho estados se declaran por nombre de variable, que es la razón por la que esta
entrega los publica como derivación con su regla. Es diseño de alta fidelidad y no una captura del
portal en ejecución.}{fig:a4-tema-paleta}

La divergencia queda declarada y no disimulada: si el archivo publica más adelante los ocho valores
oficiales, se sustituyen los derivados y se anota el cambio. Publicar un valor derivado con el
rótulo de oficial habría sido el defecto caro, porque nadie lo habría vuelto a revisar.

\subsection{El tema, medido en el portal en ejecución}

Las cifras de las tablas anteriores salen del generador. Lo que sigue sale del portal abierto en un
navegador, que es la única forma de comprobar que lo calculado es también lo pintado. Sobre las
cuatro combinaciones se midió el valor resuelto de cada propiedad personalizada, la familia
tipográfica computada y el número de elementos cuyo contenido desborda su caja.

\begin{uxtabla}{|L{3.6cm}|L{2.4cm}|L{2.4cm}|L{2.2cm}|Y|}{Medición en el navegador de las cuatro combinaciones, a 1440x900}
  \uxheadrow \thd{Combinación} & \thd{Suelo resuelto} & \thd{Suelo esperado} & \thd{Familia} & \thd{Desbordes} \\
  Corriente, claro & \texttt{F4F6F9} & \texttt{F4F6F9} & Lexend Deca & 0 \\
  Corriente, oscuro & \texttt{0A0A0C} & \texttt{0A0A0C} & Lexend Deca & 0 \\
  Institucional, claro & \texttt{FFFFFF} & \texttt{FFFFFF} & Inter & 0 \\
  Institucional, oscuro & \texttt{0B1B2B} & \texttt{0B1B2B} & Inter & 0 \\
\end{uxtabla}

Las cuatro pintan su propia paleta y ninguna hereda la del tema vecino. La familia cambia con el
tema y no con el modo, que es lo que el eje declara, y la retícula la absorbe sin un solo corte:
Inter tiene mayor altura de x que Fira Sans y el riesgo real era que una celda densa creciera.

La preferencia se comprobó además contra el HTML que devuelve el servidor, y no contra el que existe
después de hidratar: con las dos cookies puestas la etiqueta de apertura llega ya como
\texttt{<html lang="es" data-tema="institucional" data-modo="oscuro">}. Es la diferencia entre una
preferencia que se aplica y una que se corrige a la vista del lector.

\figuraux[0.86\linewidth]{figuras/a4/tema/combinacion_institucional_claro_inicio.png}{Pantalla de
inicio bajo el tema institucional en modo claro. Captura del portal en ejecución a 1440x900.}{fig:a4-tema-inst-claro}

\figuraux[0.86\linewidth]{figuras/a4/tema/combinacion_institucional_oscuro_inicio.png}{La misma
pantalla bajo el tema institucional en modo oscuro, sobre el suelo \texttt{0B1B2B}. El azul profundo
sustituye al negro puro por la regla que el sistema ya se había puesto. Captura del portal en
ejecución.}{fig:a4-tema-inst-oscuro}

\figuraux[0.86\linewidth]{figuras/a4/tema/institucional_claro_2_exploracion.png}{Catálogo temático
bajo el tema institucional. Es una de las siete pantallas capturadas bajo el tema nuevo; las quince
del resto del documento conservan el tema de omisión, que no cambió.}{fig:a4-tema-inst-catalogo}

Las capturas del tema viven en \texttt{figuras/a4/tema/} y no sustituyen a ninguna de las anteriores.
Las quince del par antes y después documentan la primera iteración bajo el tema de omisión, que esta
entrega no movió: recapturarlas habría dado imágenes más nuevas y habría destruido la comparación que
sostiene el criterio de iteración.

% ============================================================================
\uxsection{Los cinco flujos de tarea, en diseño de alta fidelidad}

Además de las siete pantallas construidas en el entorno del producto, esta actividad produjo
\textbf{cinco flujos de tarea} en el archivo de diseño: \textbf{30 pantallas} conectadas que
recorren consulta, análisis, decisión, publicación y gobierno. Los cinco comparten identidad,
retícula y componentes, porque pertenecen a un solo producto.

\textbf{Son diseño de alta fidelidad y no pantallas del producto.} La distinción no es un matiz:
cada figura de esta sección es una lámina del archivo de diseño, no una captura del portal en
ejecución, y cada pie lo dice. Las capturas del portal viven en las secciones anteriores y se
distinguen por su franja de alcance y por su barra lateral real. Mezclar las dos evidencias
convertiría un entregable verificable en una promesa.

Cada flujo se ancla a una persona de la Actividad 1 y al escenario que la Actividad 2 escribió para
ella. Los rótulos de cuenta que aparecen dentro de las láminas son datos de demostración del
archivo; la correspondencia con las personas del estudio la fija la tarea y el perfil, no el nombre
impreso en la sesión.

\begin{uxtabla}{|L{3.3cm}|L{3.3cm}|Y|C{1.5cm}|}{Los cinco flujos de tarea, con su persona, su objetivo y su número de pantallas}
  \uxheadrow \thd{Flujo} & \thd{Persona y perfil} & \thd{Objetivo de la tarea} & \thd{Pantallas} \\
  Análisis y exportación de liquidez & Diego Hernández, analista de datos & Cruzar fuentes, acotar variables y dejar una exportación reproducible & 8 \\
  Señal de riesgo para la Junta & Arturo Castañeda, directivo & Entender una señal de riesgo y llevarla explicada a una sesión de quince minutos & 5 \\
  Publicación de una versión del catálogo & Roberto Valdez, propietario del dato & Publicar una versión nueva sin romper el trabajo que depende de la anterior & 6 \\
  Administración de accesos & Mariana Ovalle Ríos, administración de la plataforma & Otorgar el acceso mínimo suficiente y dejar evidencia de la autorización & 5 \\
  Centro de trabajo integral & Ximena Solís Barrera, integración de aplicaciones & Recorrer los módulos del portal desde un solo punto de entrada, hasta las integraciones & 6 \\
\end{uxtabla}

\subsection{Flujo 1. Análisis y exportación de liquidez}

Ocho pantallas que van de la localización de la información a la confirmación del archivo
generado, pasando por el detalle de la fuente, la selección de variables, los filtros, la vista
previa y la configuración de la exportación. El paso activo, las restricciones de acceso y el
resumen de la selección permanecen visibles en cada etapa, de modo que quien exporta puede revisar
sus decisiones antes de comprometerlas.

La lámina de detalle muestra el momento en que la tarea se juega: cinco variables de veinticuatro
seleccionadas, tres de ellas marcadas como restringidas y con los permisos verificados a la vista.
Es la traducción visual del principio de que el sistema explica lo que se puede ver antes de
dejarlo pedir.

\figuraux[0.94\linewidth]{figuras/a4/secuencia_flujo_1.png}{Diseño de alta fidelidad del flujo de
análisis y exportación de liquidez: las ocho pantallas, de inicio a confirmación del
archivo.}{fig:a4-flujo-1-secuencia}

\figuraux{figuras/a4/flujo_1_variables.png}{Diseño de la pantalla de selección de variables del
flujo de análisis y exportación. Muestra el paso activo, las restricciones de acceso y el resumen
de la selección.}{fig:a4-flujo-1-detalle}

\subsection{Flujo 2. Señal de riesgo para la Junta}

Cinco pantallas compuestas para tableta, que es el equipo con el que el escenario directivo de la
Actividad 2 describe esta tarea. El recorrido va del panorama a la detección del cambio, de ahí a
su explicación y a la composición del indicador, y termina en un resumen listo para exponer. La
jerarquía privilegia el nivel de riesgo, la procedencia del dato y la evidencia que respalda la
decisión.

El flujo no se limita a mostrar una alerta: declara qué falta confirmar antes de atribuir una
causa. Es la respuesta de diseño al hallazgo de la Actividad 3, donde la tarea del indicador de
riesgo obtuvo un solo acierto de ocho contra la ruta esperada.

\figuraux[0.80\linewidth]{figuras/a4/secuencia_flujo_2.png}{Diseño de alta fidelidad del flujo de
señal de riesgo: panorama, detección, explicación, composición y resumen
ejecutivo.}{fig:a4-flujo-2-secuencia}

\figuraux[0.58\linewidth]{figuras/a4/flujo_2_riesgo.png}{Diseño de la pantalla de cambio detectado.
El estado combina magnitud, umbral, evolución de seis meses y la acción
siguiente.}{fig:a4-flujo-2-detalle}

\clearpage

\subsection{Flujo 3. Publicación de una versión del catálogo}

Seis pantallas que reparten la tarea en momentos separables: selección del elemento, documentación
del cambio, comparación entre versiones, análisis de impacto, revisión final y confirmación. La
separación no es burocracia de interfaz: es lo que permite conservar evidencia y anticipar
dependencias antes de publicar.

La lámina de detalle contrasta las dos versiones campo por campo ---definición, regla histórica,
estructura e interpretación--- y declara qué consultas hay que revisar antes del siguiente corte.
Publicar sin ese contraste es el defecto que el escenario de la Actividad 2 describe con nombre
propio.

\figuraux{figuras/a4/secuencia_flujo_3.png}{Diseño de alta fidelidad del flujo de publicación del
catálogo, desde la selección del registro hasta la confirmación de la nueva
versión.}{fig:a4-flujo-3-secuencia}

\figuraux{figuras/a4/flujo_3_versiones.png}{Diseño de la pantalla de comparación de versiones. La
interfaz contrasta definición, regla histórica, estructura e interpretación antes de
continuar.}{fig:a4-flujo-3-detalle}

\subsection{Flujo 4. Administración de accesos}

Cinco pantallas que aplican el principio de permiso mínimo suficiente: solicitud, revisión de la
justificación, definición del alcance habilitado, autorización y alta. Las confirmaciones se
reservan para las acciones que afectan a otras personas, y la pantalla final comunica vigencia,
alcance, exclusiones y evidencia disponible.

La lámina de detalle es la que sostiene el criterio: la autorización se confirma con su
justificación, quién la solicitó, quién la aprobó, su vigencia con caducidad automática, sus
exclusiones y los cuatro controles verificados. El registro queda como antecedente de cualquier
modificación posterior.

\figuraux[0.80\linewidth]{figuras/a4/secuencia_flujo_4.png}{Diseño de alta fidelidad del flujo de
administración de accesos: solicitud, revisión, permiso mínimo, autorización y
alta.}{fig:a4-flujo-4-secuencia}

\figuraux[0.58\linewidth]{figuras/a4/flujo_4_autorizacion.png}{Diseño de la pantalla de
confirmación de autorización. Conserva justificación, responsables, vigencia, alcance, exclusiones
y controles verificados.}{fig:a4-flujo-4-detalle}

\clearpage

\subsection{Flujo 5. Centro de trabajo integral}

Seis vistas que comprueban que una navegación común sostiene exploración, exportación, monitoreo de
riesgos, gobierno e integración sin perder el contexto. Más que una tarea lineal, este flujo
representa la arquitectura del producto y sus puntos de entrada, y por eso es el que verifica que
los otros cuatro caben bajo el mismo armazón.

La lámina de detalle abre con indicadores, módulos principales y actividad reciente, y su barra
lateral declara el modo de trabajo con el que se compone la vista. Es el recorrido del escenario de
integración de la Actividad 2, que termina en el módulo de credenciales y ejecuciones programadas.

\figuraux{figuras/a4/secuencia_flujo_5.png}{Diseño de alta fidelidad del flujo integral: centro de
trabajo, análisis, exportación, riesgos, gobierno e integración.}{fig:a4-flujo-5-secuencia}

\figuraux{figuras/a4/flujo_5_centro_trabajo.png}{Diseño del centro de trabajo. Reúne indicadores,
módulos principales y actividad reciente bajo una navegación común.}{fig:a4-flujo-5-detalle}

\subsection{Verificación técnica del archivo de diseño}

Antes de cerrar la entrega se revisaron la estructura, la legibilidad, los componentes y las
conexiones del archivo. La comprobación es técnica y no sustituye a una prueba de usabilidad con
participantes: su propósito es retirar los defectos evidentes antes de que el prototipo se use en
una evaluación.

\begin{uxtabla}{|L{5.1cm}|Y|}{Resultados de la verificación técnica del archivo de diseño}
  \uxheadrow \thd{Dimensión revisada} & \thd{Resultado} \\
  Pantallas y flujos & 30 pantallas organizadas en 5 flujos \\
  Interacciones & 84 nodos interactivos; 0 destinos inválidos \\
  Sistema de componentes & 12 conjuntos y 75 variantes \\
  Legibilidad & 0 textos por debajo de 12 px; 0 fuentes faltantes \\
  Integridad del contenido & 0 recortes de texto o de componente \\
  Objetivos interactivos & 0 controles por debajo de 44 por 44 px \\
  Adaptación documentada & Vistas de referencia de 1440 por 1024 y de 1024 por 1366 px \\
\end{uxtabla}

Los seis ceros de la tabla son el resultado de una revisión que sí encontró trabajo: durante ella
se sustituyeron llamadas a la acción sueltas por instancias de componente, se normalizaron los
botones primarios y se verificaron los contenedores de acciones de los diálogos. El cierre quedó
sin incidencias técnicas abiertas.

Lo que estos números no dicen es igual de importante: miden el archivo, no la experiencia. No
aportan evidencia sobre eficacia, eficiencia ni satisfacción, y esa evidencia es exactamente la que
la prueba de usabilidad de la Actividad 5 produce sobre los mismos recorridos.
